\documentclass[aspectratio=169]{beamer}

\usetheme{default}
\usecolortheme{default}
\setbeamertemplate{navigation symbols}{}
\setbeamertemplate{footline}[frame number]

\usepackage[utf8]{inputenc}
\usepackage[T1]{fontenc}
\usepackage{hyperref}
\usepackage{url}

\title{An open wiki on AO primitives}
\subtitle{ethproofs}
\author{Danny Willems}
\date{May 15, 2026}

\begin{document}

\begin{frame}
  \titlepage{}
\end{frame}

\begin{frame}{Why a wiki?}
  \begin{itemize}
    \item Started as a personal research initiative during a break to go deeper
	  into the maths and gather the whole literature.
    \item Many AO hash functions popped up in the last years:
          Poseidon, Rescue, Griffin, Anemoi, Tip5, Monolith, MiMC, \ldots
    \item If the industry continues to rely on non-binary fields
	(relatively small or large) and require more AO primitives, we will need
	to keep track of the research work and the usage of the
        primitives in the publicly available products.
    \item Security analysis, mathematics and algorithms are scattered across
	papers, PhD thesis, GitHub issues, informal notes, group discussions,
	videos, etc. No single place to look it up.
    \item Goal: a \textbf{centralised} and \textbf{live} reference for the
	cryptanalysis of AO primitives, with Lean4 and Sage code.
  \end{itemize}
\end{frame}

\begin{frame}{Current state of the wiki}
  \begin{itemize}
    \item \textbf{Primitive pages.} One per AO design, with the
          specification, security claims, and a chronological list of
          published attacks.
    \item \textbf{Three analysis axes} for each primitive:
          arithmetization cost (R1CS, AIR, PlonK-ish), CPU cost,
          cryptanalysis (algebraic, Groebner, differential,
	  statistical) with history and code in Sage.
    \item \textbf{Industrial usage map.} Which primitive is used by
          which proof system, which protocol, which deployment.
    \item \textbf{Resource gathering.} Papers and textbooks on
	  Groebner bases, multivariate solving, security analysis, etc from
	  ePrint and other resources.
  \end{itemize}
  \vspace{0.3em}
  \small \textbf{Preliminary work.} For now mostly bibliography
  gathering; the structure still needs to be reorganized.
  AI is used to go across the literature faster, including via
  Papyrus (\url{https://papyrus.badaas.eu/about}).
  It \textbf{might} contain errors, and there might be some wrong
  attributions/references. Please, notify me ASAP if you notice something
  wrong.
\end{frame}

\begin{frame}{Three directions for the project}
  \begin{itemize}
    \item \textbf{Formalization in Lean.} Mechanize all definitions,
          lemmas, and theorems used to analyse AO primitives.
          Verify the pen-and-paper proofs. Track the open questions
          explicitly.
    \item \textbf{Reproducible attacks in Sage.} A centralized
          open-source repository of executable attack scripts for
          every published attack, in the spirit of \texttt{ECC\_Attacks}
          for elliptic curves.
    \item \textbf{New cryptanalysis.} Independently, following from the first two: explore underused directions,
	    e.g.\ geometry-based attacks, new Groebner basis
	    techniques, extend the mathematics used in the cryptanalysis.
  \end{itemize}
\end{frame}

\begin{frame}{Status and how to contribute}
  \begin{itemize}
    \item Open source, open to PRs: corrections, new primitives,
          new attacks, missing references.
    \item Wiki: \url{https://badaas.be/ao-security}
    \item Repo: \url{https://github.com/BaDaaS/ao-security}
  \end{itemize}
  \vspace{1em}
  \centering
  \large Thank you. Questions?
\end{frame}

\end{document}
